% Figure from MATH 590 notes, section 7. Compile with the notes preamble.
\begin{tikzpicture}[use Hobby shortcut]
  \def\Ablob{(0,0) .. (1.5,-0.4) .. (2.8,0.3) .. (2.7,1.6) .. (1.3,2.0) .. (0.1,1.3)}
  \path[fill=figB!10, name path=A, closed] \Ablob;
  % part of the boundary belongs to A (solid), part does not (dashed)
  \begin{scope}
    \clip (-1,-1) rectangle (1.6,3);
    \path[draw=figB, thick, closed] \Ablob;
  \end{scope}
  \begin{scope}
    \clip (1.6,-1) rectangle (4,3);
    \path[draw=figB, thick, dashed, closed] \Ablob;
  \end{scope}
  \node[font=\small, figB] at (0.45,0.35) {$A$};
  % interior point p: a neighborhood inside A
  \coordinate (p) at (1.0,1.05);
  \draw[figB, thick, fill=figB!20] (p) circle (0.32);
  \fill (p) circle (1.3pt);
  \node[font=\scriptsize, right=0pt] at (p) {$p$};
  % boundary point x (not in A): every neighborhood meets A and its complement
  \path[name path=h] (2.0,0.95) -- (3.5,0.95);
  \path[name intersections={of=A and h, by=x}];
  \begin{scope}
    \path[clip, closed] \Ablob;
    \fill[figR!22] (x) circle (0.38);
  \end{scope}
  \draw[figR, thick] (x) circle (0.38);
  \draw[figR, thin] (x) circle (0.2);
  \fill[figR] (x) circle (1.3pt);
  \node[font=\scriptsize, figR, right] at ($(x)+(0.38,0)$) {$x$};
  % exterior point y: a neighborhood U missing A
  \coordinate (y) at (3.7,2.0);
  \draw[black!55, thick, fill=black!5] (y) circle (0.4);
  \fill (y) circle (1.3pt);
  \node[font=\scriptsize, right=0pt] at (y) {$y$};
  \node[font=\scriptsize, black!60] at ($(y)+(0.12,-0.58)$) {$U$};
\end{tikzpicture}
